% GENERATED FILE. Edit canonical lessons, then run npm run generate:books.
% canonical-source-hash: fnv1a64:7da0dafd72fed04e
% canonical-lessons: MW-C240-pakno, MW-C240-lunno, MW-C240-khalo, MW-C240-bali, MW-C240-dher, MW-R240-review-235-240, MW-R240-review-final-240

\chapter{Harvest}
\label{ch:mw-harvest}
\hlchaptermodality{240}

\begin{chapteropening}
\textbf{By the end of this chapter:} \emph{I can say the millet ripens and I will reap it.}

Complete the last lesson of chapter 240: I can say the millet ripens and I will reap it.
\end{chapteropening}

\section[\texorpdfstring{\mw{पाकणो} (pākṇo)}{पाकणो (pākṇo)}]{\mw{पाकणो} (pākṇo) — to ripen}

\label{lesson:MW-C240-pakno}



\begin{quote}
Before the new one: say the Marwadi for to grow, then the Marwadi for a new shoot.
\end{quote}

\subsection*{You'll want to know: \mw{पाकणो}}
\textbf{\mw{पाकणो}} — \emph{pākṇo} — \textquotedblleft{}to ripen\textquotedblright{}.

\begin{quote}
\mw{काती} \mw{मांय} \mw{बाजरो} \mw{पाकै} \mw{है।}
\end{quote}

\begin{quote}
\emph{kātī mā̃y bājro pākai hai.}
\end{quote}

\begin{quote}
The millet ripens in autumn.
\end{quote}

Past (it …): \mw{वो} \mw{बाजरो} \mw{पाक्यो}, \emph{vo bājro pākyo}.

\begin{cousinweb}
\textbf{\mw{पाकणो}} (\emph{pākṇo}) is inherited from Sanskrit \emph{pakva}, \textquotedblleft{}cooked, ripe\textquotedblright{}; CDIAL lists the Old Marwari verb \emph{pākaï}, \textquotedblleft{}becomes ripe\textquotedblright{} (CDIAL 7621). \textbf{\mw{पक्को}}, settled, comes from the same word.
\end{cousinweb}

\begin{grammarlens}[title={What you've built}]
One more word for this chapter.
\end{grammarlens}

\subsection*{Your turn}
\begin{itemize}
\raggedright
  \item \emph{Say it:} \emph{pākṇo}
  \item \emph{Say it:} \emph{pākṇo}, once more
  \item \emph{Say it:} \emph{kātī mā̃y bājro pākai hai.} — The millet ripens in autumn.
  \item \emph{Recall:} say the Marwadi for to grow, then the Marwadi for a new shoot, then say \emph{pākṇo} again
\end{itemize}

\begin{tcolorbox}[breakable,colback=teal!4,colframe=teal!35!black,title={Before you move on}]
What does \mw{पाकणो} mean? (\textquotedblleft{}To ripen\textquotedblright{}.) Say it once more.
\end{tcolorbox}
\section[\texorpdfstring{\mw{लूणणो} (lūṇṇo)}{लूणणो (lūṇṇo)}]{\mw{लूणणो} (lūṇṇo) — to reap}

\label{lesson:MW-C240-lunno}



\begin{quote}
Before the new one: say the Marwadi for a new shoot, then the Marwadi for to ripen.
\end{quote}

\subsection*{You'll want to know: \mw{लूणणो}}
\textbf{\mw{लूणणो}} — \emph{lūṇṇo} — \textquotedblleft{}to reap\textquotedblright{}.

\begin{quote}
\mw{हूं} \mw{काल} \mw{बाजरो} \mw{लूणसूं।}
\end{quote}

\begin{quote}
\emph{hū̃ kāl bājro lūṇsū̃.}
\end{quote}

\begin{quote}
Tomorrow I'll reap the millet.
\end{quote}

Past (I …): \mw{म्हैं} \mw{बाजरो} \mw{लूण्यो}, \emph{mhaĩ bājro lūṇyo}.

\begin{cousinweb}
\textbf{\mw{लूणणो}} (\emph{lūṇṇo}) is inherited from Sanskrit \emph{lunāti}, \textquotedblleft{}cuts, reaps\textquotedblright{}; CDIAL lists Hindi \emph{lunnā}, \textquotedblleft{}to reap\textquotedblright{}. CDIAL gives no Marwari form in this entry (CDIAL 11082).
\end{cousinweb}

\begin{grammarlens}[title={What you've built}]
One more word for this chapter.
\end{grammarlens}

\subsection*{Your turn}
\begin{itemize}
\raggedright
  \item \emph{Say it:} \emph{lūṇṇo}
  \item \emph{Say it:} \emph{lūṇṇo}, once more
  \item \emph{Say it:} \emph{hū̃ kāl bājro lūṇsū̃.} — Tomorrow I'll reap the millet.
  \item \emph{Recall:} say the Marwadi for a new shoot, then the Marwadi for to ripen, then say \emph{lūṇṇo} again
\end{itemize}

\begin{tcolorbox}[breakable,colback=teal!4,colframe=teal!35!black,title={Before you move on}]
What does \mw{लूणणो} mean? (\textquotedblleft{}To reap\textquotedblright{}.) Say it once more.
\end{tcolorbox}
\section[\texorpdfstring{\mw{खळो} (khaḷo)}{खळो (khaḷo)}]{\mw{खळो} (khaḷo) — a threshing floor}

\label{lesson:MW-C240-khalo}



\begin{quote}
Before the new one: say the Marwadi for to ripen, then the Marwadi for to reap.
\end{quote}

\subsection*{You'll want to know: \mw{खळो}}
\textbf{\mw{खळो}} — \emph{khaḷo} — \textquotedblleft{}a threshing floor\textquotedblright{}.

\begin{quote}
\mw{खळो} \mw{खेत} \mw{कनै} \mw{है।}
\end{quote}

\begin{quote}
\emph{khaḷo khet kanai hai.}
\end{quote}

\begin{quote}
The threshing floor is near the field.
\end{quote}

A masculine noun.

\begin{cousinweb}
\textbf{\mw{खळो}} (\emph{khaḷo}) is inherited from Sanskrit \emph{khala}, \textquotedblleft{}a threshing floor\textquotedblright{}; CDIAL lists Gujarati \emph{khaḷũ}, \textquotedblleft{}a threshing floor\textquotedblright{}. CDIAL gives no Marwari form in this entry (CDIAL 3834).
\end{cousinweb}

\begin{grammarlens}[title={What you've built}]
One more word for this chapter.
\end{grammarlens}

\subsection*{Your turn}
\begin{itemize}
\raggedright
  \item \emph{Say it:} \emph{khaḷo}
  \item \emph{Say it:} \emph{khaḷo}, once more
  \item \emph{Say it:} \emph{khaḷo khet kanai hai.} — The threshing floor is near the field.
  \item \emph{Recall:} say the Marwadi for to ripen, then the Marwadi for to reap, then say \emph{khaḷo} again
\end{itemize}

\begin{tcolorbox}[breakable,colback=teal!4,colframe=teal!35!black,title={Before you move on}]
What does \mw{खळो} mean? (\textquotedblleft{}A threshing floor\textquotedblright{}.) Say it once more.
\end{tcolorbox}
\section[\texorpdfstring{\mw{बाली} (bālī)}{बाली (bālī)}]{\mw{बाली} (bālī) — an ear of grain}

\label{lesson:MW-C240-bali}



\begin{quote}
Before the new one: say the Marwadi for to reap, then the Marwadi for a threshing floor.
\end{quote}

\subsection*{You'll want to know: \mw{बाली}}
\textbf{\mw{बाली}} — \emph{bālī} — \textquotedblleft{}an ear of grain\textquotedblright{}.

\begin{quote}
\mw{बाली} \mw{पीळी} \mw{है।}
\end{quote}

\begin{quote}
\emph{bālī pīḷī hai.}
\end{quote}

\begin{quote}
The ear of grain is yellow.
\end{quote}

A feminine noun.

\begin{cousinweb}
\textbf{\mw{बाली}} (\emph{bālī}) is inherited from Sanskrit \emph{valla}, \textquotedblleft{}a kind of wheat\textquotedblright{}; CDIAL lists Old Marwari \emph{bālī}, \textquotedblleft{}the standing crop\textquotedblright{} (CDIAL 11425). In Marwadi it names the ear at the top of the stalk.
\end{cousinweb}

\begin{grammarlens}[title={What you've built}]
One more word for this chapter.
\end{grammarlens}

\subsection*{Your turn}
\begin{itemize}
\raggedright
  \item \emph{Say it:} \emph{bālī}
  \item \emph{Say it:} \emph{bālī}, once more
  \item \emph{Say it:} \emph{bālī pīḷī hai.} — The ear of grain is yellow.
  \item \emph{Recall:} say the Marwadi for to reap, then the Marwadi for a threshing floor, then say \emph{bālī} again
\end{itemize}

\begin{tcolorbox}[breakable,colback=teal!4,colframe=teal!35!black,title={Before you move on}]
What does \mw{बाली} mean? (\textquotedblleft{}An ear of grain\textquotedblright{}.) Say it once more.
\end{tcolorbox}
\section[\texorpdfstring{\mw{ढेर} (ḍher)}{ढेर (ḍher)}]{\mw{ढेर} (ḍher) — a heap; a pile}

\label{lesson:MW-C240-dher}



\begin{quote}
Before the new one: say the Marwadi for a threshing floor, then the Marwadi for an ear of grain.
\end{quote}

\subsection*{You'll want to know: \mw{ढेर}}
\textbf{\mw{ढेर}} — \emph{ḍher} — \textquotedblleft{}a heap; a pile\textquotedblright{}.

\begin{quote}
\mw{खेत} \mw{मांय} \mw{मोटो} \mw{ढेर} \mw{है।}
\end{quote}

\begin{quote}
\emph{khet mā̃y moṭo ḍher hai.}
\end{quote}

\begin{quote}
There's a big heap in the field.
\end{quote}

A masculine noun.

\begin{cousinweb}
\textbf{\mw{ढेर}} (\emph{ḍher}) is inherited from a reconstructed \emph{*ḍhēra}, \textquotedblleft{}a heap\textquotedblright{}; CDIAL lists Hindi and Gujarati \emph{ḍher}. CDIAL gives no Marwari form in this entry (CDIAL 5599).
\end{cousinweb}

\begin{grammarlens}[title={What you've built}]
One more word for this chapter.
\end{grammarlens}

\subsection*{Your turn}
\begin{itemize}
\raggedright
  \item \emph{Say it:} \emph{ḍher}
  \item \emph{Say it:} \emph{ḍher}, once more
  \item \emph{Say it:} \emph{khet mā̃y moṭo ḍher hai.} — There's a big heap in the field.
  \item \emph{Recall:} say the Marwadi for a threshing floor, then the Marwadi for an ear of grain, then say \emph{ḍher} again
\end{itemize}

\begin{tcolorbox}[breakable,colback=teal!4,colframe=teal!35!black,title={Before you move on}]
What does \mw{ढेर} mean? (\textquotedblleft{}A heap; a pile\textquotedblright{}.) Say it once more.
\end{tcolorbox}
\section[(dialogue)]{Looking back: from The Cattle Yard to Harvest}

\label{lesson:MW-R240-review-235-240}



\begin{quote}
Say the Marwadi for a cattle yard and fodder.
\end{quote}

\subsection*{Your turn}
Say these aloud:

\begin{itemize}
\raggedright
  \item a cattle yard, fodder, cow dung, to rear (animals), chaff
  \item a donkey, a rooster, wool, a horn, a hoof
  \item a partridge, a wing, to peck, an egg, a frog
  \item to sow, to dig, to plant (a seedling or a tree), a ploughman, a seed
  \item sorghum, mustard, cotton, to grow, a new shoot
  \item to ripen, to reap, a threshing floor, an ear of grain, a heap
\end{itemize}

\begin{tcolorbox}[breakable,colback=teal!4,colframe=teal!35!black,title={Before you move on}]
A cattle yard? (\textbf{\mw{गुवाड़ो}}, \emph{guvāṛo}.) Fodder? (\textbf{\mw{चारो}}, \emph{chāro}.) Cow dung? (\textbf{\mw{गोबर}}, \emph{gobar}.) To rear (animals)? (\textbf{\mw{पाळणो}}, \emph{pāḷṇo}.) Chaff? (\textbf{\mw{भूसो}}, \emph{bhūso}.) A donkey? (\textbf{\mw{गधो}}, \emph{gadho}.) A rooster? (\textbf{\mw{कूकड़ो}}, \emph{kūkṛo}.) Wool? (\textbf{\mw{ऊन}}, \emph{ūn}.) A horn? (\textbf{\mw{सींग}}, \emph{sī̃g}.) A hoof? (\textbf{\mw{खुर}}, \emph{khur}.) A partridge? (\textbf{\mw{तीतर}}, \emph{tītar}.) A wing? (\textbf{\mw{पांख}}, \emph{pā̃kh}.) To peck? (\textbf{\mw{चुगणो}}, \emph{chugṇo}.) An egg? (\textbf{\mw{इंडो}}, \emph{iṇḍo}.) A frog? (\textbf{\mw{मींडको}}, \emph{mī̃ḍko}.) To sow? (\textbf{\mw{बोवणो}}, \emph{bovṇo}.) To dig? (\textbf{\mw{खोदणो}}, \emph{khodṇo}.) To plant (a seedling or a tree)? (\textbf{\mw{रोपणो}}, \emph{ropṇo}.) A ploughman? (\textbf{\mw{हाळी}}, \emph{hāḷī}.) A seed? (\textbf{\mw{बीज}}, \emph{bīj}.) Sorghum? (\textbf{\mw{जुवार}}, \emph{juvār}.) Mustard? (\textbf{\mw{सरसूं}}, \emph{sarsū̃}.) Cotton? (\textbf{\mw{कपास}}, \emph{kapās}.) To grow? (\textbf{\mw{बधणो}}, \emph{badhṇo}.) A new shoot? (\textbf{\mw{कूंपळ}}, \emph{kū̃paḷ}.) To ripen? (\textbf{\mw{पाकणो}}, \emph{pākṇo}.) To reap? (\textbf{\mw{लूणणो}}, \emph{lūṇṇo}.) A threshing floor? (\textbf{\mw{खळो}}, \emph{khaḷo}.) An ear of grain? (\textbf{\mw{बाली}}, \emph{bālī}.) A heap? (\textbf{\mw{ढेर}}, \emph{ḍher}.)
\end{tcolorbox}
\section[(dialogue)]{Harvest — once more}

\label{lesson:MW-R240-review-final-240}



\begin{quote}
Say the Marwadi for to ripen and to reap.
\end{quote}

\subsection*{Your turn}
\begin{itemize}
\raggedright
  \item \emph{Say it:} to ripen, to reap, a threshing floor, an ear of grain, a heap
\end{itemize}

\begin{tcolorbox}[breakable,colback=teal!4,colframe=teal!35!black,title={Before you move on}]
To ripen? (\textbf{\mw{पाकणो}}, \emph{pākṇo}.) To reap? (\textbf{\mw{लूणणो}}, \emph{lūṇṇo}.) A threshing floor? (\textbf{\mw{खळो}}, \emph{khaḷo}.) An ear of grain? (\textbf{\mw{बाली}}, \emph{bālī}.) A heap? (\textbf{\mw{ढेर}}, \emph{ḍher}.)
\end{tcolorbox}
